\documentclass[11pt,a4paper]{article}

% --- Pacotes Fundamentais ---
\usepackage[utf8]{inputenc}
\usepackage[brazil,english]{babel}
\usepackage{amsmath,amssymb,amsfonts}
\usepackage{graphicx}
\usepackage{booktabs}
\usepackage{hyperref}
\usepackage{cite}
\usepackage{geometry}
\geometry{margin=2.5cm}

% --- Informações do Artigo ---
\title{\textbf{Investigação e Implementação de Modelos Neurais em Tópicos Avançados de Inteligência Artificial}}

\author{
    Aluno: [Seu Nome Aqui] \\
    \textit{Programa de Pós-Graduação / Graduação em Ciência da Computação} \\
    \textit{Disciplina: Tópicos Avançados em Inteligência Artificial}
}

\date{\today}

\begin{document}

\selectlanguage{brazil}

\maketitle

\begin{abstract}
Este trabalho apresenta o estudo, projeto e implementação de arquiteturas neurais avançadas no âmbito da disciplina de Tópicos Avançados em Inteligência Artificial. Através de uma abordagem baseada em grafo explícito e experimentação empírica, investigam-se os impactos de hiperparâmetros, dinâmicas de treinamento e convergência de gradientes. Os resultados obtidos demonstram a eficácia dos métodos propostos, validando a integração entre o ambiente de codificação e a análise técnica formal.
\end{abstract}

\vspace{0.5cm}
\noindent\textbf{Palavras-chave:} Inteligência Artificial, Redes Neurais Artificiais, Multi-Layer Perceptron, Otimização de Hiperparâmetros.

\vspace{0.8cm}

% --- Importação Modular das Seções ---
\section{Introdução}
\label{sec:introducao}

O avanço contínuo das técnicas de Inteligência Artificial (IA) e Aprendizado Profundo (\textit{Deep Learning}) tem impulsionado transformações significativas em diversas áreas do conhecimento \cite{goodfellow2016deep}. No entanto, a compreensão detalhada das dinâmicas internas de treinamento, fluxos de propagação direta (\textit{forward propagation}) e retropropagação de gradientes (\textit{backpropagation}) permanece fundamental para o desenvolvimento e depuração de novos modelos.

Este artigo propõe e avalia a implementação de modelos neurais sob a perspectiva de grafos explícitos, permitindo a inspeção detalhada de estados internos e a visualização dinâmica da arquitetura. O objetivo central é fornecer uma base empírica e analítica sólida para as atividades desenvolvidas na disciplina de Tópicos Avançados em Inteligência Artificial.

\section{Trabalhos Relacionados}
\label{sec:trabalhos_relacionados}

Os fundamentos do treinamento de redes neurais multicamadas através da descida de gradiente e retropropagação foram consolidados no trabalho seminal de Rumelhart et al. \cite{rumelhart1986learning}. Desde então, diversos arcabouços computacionais foram propostos para simplificar o cálculo automático de gradientes e a construção de topologias complexas.

Na literatura recente, destacam-se abordagens voltadas à interpretabilidade e representação explícita de grafos de computação, as quais auxiliam pesquisadores e estudantes a diagnosticar fenômenos como o desvanecimento de gradientes (\textit{vanishing gradients}) e a saturação de funções de ativação não lineares.

\section{Metodologia}
\label{sec:metodologia}

A metodologia adotada divide-se em duas frentes complementares: formulação teórica do modelo neural e implementação em código modular.

\subsection{Modelo Matemático}
Dado um vetor de entrada $\mathbf{x} \in \mathbb{R}^{d}$, a saída de um neurônio $j$ em uma camada $l$ é expressa por:
\begin{equation}
    z_j^{(l)} = \sum_{i} w_{ij}^{(l)} a_i^{(l-1)} + b_j^{(l)}
\end{equation}
\begin{equation}
    a_j^{(l)} = \sigma\left(z_j^{(l)}\right)
\end{equation}
onde $w_{ij}^{(l)}$ denota o peso sináptico da conexão, $b_j^{(l)}$ representa o termo de viés (\textit{bias}) e $\sigma(\cdot)$ é a função de ativação não linear aplicada elemento a elemento.

\subsection{Arquitetura em Grafo Explícito}
Diferentemente de implementações baseadas exclusivamente em operações matriciais opacas, a arquitetura foi concebida como um grafo direcionado $G = (V, E)$, onde os vértices $V$ correspondem aos neurônios individuais e as arestas $E$ armazenam os pesos sinápticos e os respectivos gradientes acumulados em cada passo de otimização.

\section{Resultados e Discussão}
\label{sec:experimentos}

Esta seção detalha os experimentos empíricos conduzidos com a arquitetura proposta, avaliando convergência de função de perda, taxa de acerto (\textit{accuracy}) e estabilidade numérica.

\subsection{Ambiente Experimental}
Os experimentos foram executados utilizando o ambiente integrado em Python, com dados normalizados e divisão entre conjuntos de treino e validação.

\begin{table}[ht]
\centering
\caption{Resumo dos hiperparâmetros utilizados nos ensaios experimentais.}
\label{tab:hiperparametros}
\begin{tabular}{@{}ll@{}}
\toprule
\textbf{Parâmetro} & \textbf{Valor Adotado} \\ \midrule
Taxa de Aprendizado ($\eta$) & 0.01 -- 0.1 \\
Otimizador & Descida do Gradiente Estocástico (SGD) \\
Função de Custo & Erro Quadrático Médio (MSE) / Entropia Cruzada \\
Épocas de Treinamento & 500 \\ \bottomrule
\end{tabular}
\end{table}

Os resultados preliminares evidenciam a convergência monótona da perda e a consistência dos gradientes calculados através da estrutura em grafo.

\section{Conclusão e Trabalhos Futuros}
\label{sec:conclusao}

Este trabalho apresentou uma abordagem integrada para o estudo e implementação de modelos neurais em inteligência artificial. A representação explícita como grafo facilitou a depuração algorítmica e a inspeção dos gradientes durante a etapa de treinamento.

Como trabalhos futuros, planeja-se a extensão para arquiteturas baseadas em mecanismos de atenção \cite{vaswani2017attention}, bem como a automatização da exportação de métricas experimentais diretamente para o relatório acadêmico.


% --- Referências Bibliográficas ---
\bibliographystyle{plain}
\bibliography{references}

\end{document}
